\chapter{kb-aads: The Agent Swarm (Optional)}
\label{ch:aads}
\cmd{kb-aads} is a Python/Ray multi-agent system that consumes \cmd{kbd}'s data and coordinates response. It has
\textbf{no compiled build}. The package installs its sources under \file{/usr/lib/kb/aads/} and a launcher,
\cmd{kbagents}; it does \textbf{not} install Ray or PyTorch, which are several gigabytes and CPU-specific
(the development virtualenv is 5.4~GB). The unit is installed but not enabled.

\begin{kbnote}[STATUS: RUN AND VERIFIED, WITHOUT SYSTEMD]
The packaged launcher and sources were run on the Ubuntu 24.04 VM with that machine's existing virtualenv
(Ray 2.56): Ray started a local instance, the status server listened on \texttt{127.0.0.1:8601}, and
\texttt{GET /agents} returned nine agents (executor, judge, two hunters, two patrollers, and more), all
\texttt{active} with \texttt{error\_count} 0; the containment agent loaded the shipped checkpoint. The same launcher
also reached the swarm banner on Arch with a scratch Python 3.11 venv (Ray 2.59, CPU torch), where Ray's own
startup then stalled on a networking problem of that machine. \textbf{Not tested:} running it as
\unit{kbagents.service} under \cmd{systemd}, with a live \cmd{kbd} behind it.
\end{kbnote}

\section{Create the environment}
\resetsteps
\step{Install Python tooling}
\ifarch
\begin{lstlisting}[style=console]
$ sudo pacman -S --needed python python-pip
\end{lstlisting}
\else
\begin{lstlisting}[style=console]
$ sudo apt install -y python3 python3-venv python3-pip
\end{lstlisting}
\fi
\step{Build the venv where the launcher looks for it}
\begin{lstlisting}[style=console]
$ sudo python3 -m venv /opt/kb/venv
$ sudo /opt/kb/venv/bin/pip install -r /usr/lib/kb/aads/requirements.txt
\end{lstlisting}
Requirements: \texttt{ray[rllib,default]>=2.9.0}, \texttt{grpcio>=1.50.0}, \texttt{grpcio-tools}, \texttt{pyyaml}, \texttt{gymnasium},
\texttt{torch>=2.0.0}, plus \texttt{pytest} for tests. Expect a long download and several GB of disk.
\step{Check the launcher is satisfied}
\begin{lstlisting}[style=console]
$ kbagents --help
usage: main.py [-h] [--config CONFIG]
\end{lstlisting}
Before the venv exists, \cmd{kbagents} exits 1 and prints exactly the \texttt{pip install} command to run.

\section{Run it}
To try it without installing the package, point the launcher at an unpacked tree and an existing venv:
\begin{lstlisting}[style=console]
$ KB_AADS_DIR=/path/to/usr/lib/kb/aads KB_AADS_PYTHON=/path/to/venv/bin/python3 \
  KB_AADS_CONFIG=/path/to/agents.yaml RAY_AUTH_MODE=disabled kbagents
$ curl -s http://127.0.0.1:8601/agents
\end{lstlisting}
The status server answers on \texttt{/agents} (the bare \texttt{/} returns nothing). \texttt{RAY\_AUTH\_MODE=disabled} stops newer Ray
versions from generating a token file in your home directory; leave it unset in the service.
To run it for real:
\begin{lstlisting}[style=console]
$ sudo systemctl enable --now kbagents
$ systemctl status kbagents
$ curl -s http://127.0.0.1:8601/agents   # swarm status server
\end{lstlisting}
\texttt{kbagents.service} runs \texttt{/usr/bin/kbagents}, which \texttt{cd}s to \file{/usr/lib/kb/aads} and executes
\texttt{main.py --config /etc/kb/agents.yaml}. It \texttt{Requires=kbd.service}; it deliberately does
\textbf{not} require the maintenance-held \cmd{kb-checker}.

\section{How \cmd{kbd} finds it}
\cmd{kbd} proxies the swarm's status for the dashboard's \texttt{/api/agents} route. The status server binds
\texttt{127.0.0.1:8601} by default (\texttt{KB\_AADS\_STATUS\_HOST}, \texttt{KB\_AADS\_STATUS\_PORT}); \cmd{kbd}'s
\texttt{KB\_AADS\_API\_ADDR} must point at the same address if you change it.

\begin{kbnote}[KNOWN GAP]
The event contract notes that nothing in \cmd{kb-aads} yet \emph{consumes} \cmd{kbd}'s \texttt{process\_state} /
\texttt{zone\_transition} events. The swarm is wired and runs, but event-driven behaviour is still being built
out. Do not present it as autonomous response.
\end{kbnote}

\section{Hardening in the unit}
\texttt{NoNewPrivileges}, \texttt{PrivateTmp}, \texttt{ProtectHome}, \texttt{MemoryDenyWriteExecute},
\texttt{ProtectSystem=strict}, an empty capability set, and read-only access to \file{/etc/kb}, \file{/usr/lib/kb},
\file{/opt/kb}; only \file{/run/kb} is writable. If Ray or PyTorch turn out to need JIT memory
(\texttt{MemoryDenyWriteExecute} blocks it), relax that single directive in a drop-in and document why.
